\section{Introducción}

El presente documento tiene como objetivo aplicar los principios de la Teoría General de Sistemas (GST) para identificar y estructurar los elementos fundamentales de un \textbf{Sistema Computarizado de Gestión de Mantenimiento (CMMS)} o un sistema de \textbf{Gestión de Activos Empresariales (EAM)}, basado en un \textbf{Gemelo Digital Interactivo} para plantas industriales. Este ejercicio constituye la primera fase de análisis para el software a construir, cuyo propósito es optimizar la planificación, ejecución y análisis del mantenimiento de activos, centralizando la información y digitalizando los flujos de trabajo en campo.

Siguiendo la metodología de Análisis de Procesos de Negocio (BPA) propuesta en el material de estudio, la cual se fundamenta en la GST, este informe descompone el sistema en sus componentes clave: actores, procesos, entradas, salidas y relaciones \parencite{sena:fundamentos2024, ossa:2016}. Este informe busca establecer una base sólida y bien definida para las futuras fases de diseño y desarrollo, asegurando que la solución propuesta responda de manera coherente y completa a las necesidades de un entorno industrial moderno.
